\documentclass{article}
\usepackage{amsmath,amssymb}
\usepackage{xurl}
\usepackage[colorlinks=true,linkcolor=blue,citecolor=blue,urlcolor=blue]{hyperref}
% Shared commands for notes in docs/paper-gaps/.

\DeclareUnicodeCharacter{2080}{\ifmmode{}_{0}\else\textsubscript{0}\fi}
\DeclareUnicodeCharacter{2081}{\ifmmode{}_{1}\else\textsubscript{1}\fi}
\DeclareUnicodeCharacter{2082}{\ifmmode{}_{2}\else\textsubscript{2}\fi}
\DeclareUnicodeCharacter{2099}{\ifmmode{}_{n}\else\textsubscript{n}\fi}

\newcommand{\Lean}{\textsc{Lean}}
\ExplSyntaxOn
\NewDocumentCommand{\leanid}{m}{
  \ifmmode
    \textsf{\detokenize{#1}}
  \else
    \group_begin:
    \urlstyle{sf}
    \tl_set:Nn \l_tmpa_tl {#1}
    \tl_replace_all:Nnn \l_tmpa_tl {\_} {_}
    \exp_args:NV \path \l_tmpa_tl
    \group_end:
  \fi
}
\ExplSyntaxOff
\providecommand{\tr}{\operatorname{tr}}
\newcommand{\tnleanIssueBase}{https://github.com/LionSR/TNLean/issues/}
\newcommand{\tnleanPrBase}{https://github.com/LionSR/TNLean/pull/}
\newcommand{\tnleanDocBase}{https://sirui-lu.com/TNLean/docs/}
\newcommand{\ghissue}[1]{\href{\tnleanIssueBase#1}{GitHub issue \##1}}
\newcommand{\ghpr}[1]{\href{\tnleanPrBase#1}{PR \##1}}
\newcommand{\doclink}[1]{\href{\tnleanDocBase#1.html}{\path{#1}}}

% Dirac notation and the identity operator, as in blueprint/src/macros/common.tex.
% \providecommand keeps note-local definitions authoritative.
\usepackage{dsfont}
\providecommand{\ket}[1]{|#1\rangle}
\providecommand{\bra}[1]{\langle #1|}
\providecommand{\braket}[2]{\langle #1 | #2 \rangle}
\providecommand{\ketbra}[2]{|#1\rangle\!\langle #2|}
\providecommand{\Id}{\mathds{1}}
\providecommand{\one}{\mathds{1}}
\providecommand{\C}{\mathbb{C}}
\providecommand{\MN}[1]{M_{#1}(\C)}
\providecommand{\MD}{M_D(\C)}

% Verdict marker: \gapnote{<kind>}{<status>}. Kind is the policy
% classification (clarification, local-correction, scope-restriction,
% unfaithful, false-source, open-gap); status is open, wip (elimination
% underway), resolved, or historical. Machine-read by the site index;
% prints a verdict line.
\newcommand{\gapnote}[2]{\par\noindent\textbf{Verdict:} #1 (#2).\par}

\setlength{\emergencystretch}{2em}
\title{The Error of the Matrix Product State in Hastings' Area Law}
\author{}
\date{October 22, 2026}
\begin{document}
\maketitle
\gapnote{clarification}{resolved}

\section{The source statement}

Write $D>1$ for the local dimension, $\xi'>0$ for the decay length of the
approximate factorization, $k_0=\exp(2S)/2$ for the bond dimension set by the
entropy $S$ of a cut, and $c=4C_1^2$ for the constant of the Schmidt-tail
constraint, as in the note \path{hastings07_schmidt_tail_constant.tex}.
The section on matrix product states of~\cite{Hastings2007AreaLaw} concludes
that the error of a matrix product state of bond dimension $k'$ on a chain of
length $N$ scales as
\begin{align}
    N\,(k'/k_0)^{1/\xi'\ln(D)}.
    \label{eq:hastings07_mps_error_source}
\end{align}
It does not name the error measure.\footnote{\path{References/0705.2024/law16.tex},
    lines~553--557.} The exponent is printed without its minus sign; the bound
is meant to decay in $k'$, and the corrected exponent $-1/(\xi'\ln D)$ is
recorded in the separate note \path{hastings07_schmidt_tail_constant.tex}. The
rest of this note uses the corrected exponent.

\section{The reading}

Let $\Psi_0$ be a unit vector on an open chain of $N$ sites, and let
$\Phi$ be the vector obtained by truncating every internal cut onto its $k'$
leading Schmidt vectors. No normalization is performed. Let $\tau_j(k')$
be the sum of the squared Schmidt coefficients beyond the $k'$ largest at the cut
$(j,j+1)$. For $N<2$, the truncation is the identity and $\Phi=\Psi_0$.
For $N\ge2$, the squared Hilbert-space error of sequential truncation satisfies
\begin{align}
    \lVert\Phi-\Psi_0\rVert^2
     & \le\sum_{j=1}^{N-1}\tau_j(k'),
    \label{eq:hastings07_mps_error_sq}
\end{align}
since at each step $\Psi_0-\Pi\Phi'=\Pi(\Psi_0-\Phi')+(1-\Pi)\Psi_0$ is an
orthogonal decomposition. Verstraete and
Cirac~\cite{VerstraeteCirac2006Faithfully} state
(\ref{eq:hastings07_mps_error_sq}) with a factor $2$ on the right. Together
with the polynomial decay
$\tau_j(k')\le c\,e^{2/\xi'}(k'/k_0)^{-1/(\xi'\ln D)}$ of the tail, it gives
(\ref{eq:hastings07_mps_error_source}) for the squared norm:
\begin{align}
    \lVert\Phi-\Psi_0\rVert^2
     & \le(N-1)\,c\,e^{2/\xi'}\,(k'/k_0)^{-1/(\xi'\ln D)}.
    \label{eq:hastings07_mps_error_sq_final}
\end{align}
The error in (\ref{eq:hastings07_mps_error_source}) is therefore read as the
squared norm, which is the measure for which the stated rate
holds.\footnote{\leanid{TNLean.MPS.AreaLaw.Hastings.isOpenMPS_norm_sub_sq_le_of_schmidtTail_geometric}
    in \doclink{TNLean/MPS/AreaLaw/Hastings/SchmidtTailMPS};
    \leanid{TNLean.MPS.AreaLaw.norm_sub_sequentialTruncation_sq_le} in
    \doclink{TNLean/MPS/AreaLaw/SequentialTruncation}.}
This is a vector-norm estimate, not the trace norm of a difference of
density operators. It does not assert the same constant for a normalized
approximant. If $\eta=\lVert\Phi-\Psi_0\rVert<1$, then
$\lVert\Phi\rVert\ge1-\eta>0$, and the normalized vector
$\widehat\Phi=\Phi/\lVert\Phi\rVert$ has the same bond bound and satisfies
$\lVert\widehat\Phi-\Psi_0\rVert\le2\eta$ by the triangle and reverse
triangle inequalities. Its squared error is consequently at most four times
that of the unnormalized vector.

The norm of the unnormalized error itself is of order
$\sqrt N\,(k'/k_0)^{-1/(2\xi'\ln D)}$.

\bibliographystyle{alpha}
\bibliography{references}

\end{document}
